% The integrated commentary of the concise study, beginning on physical page
% 5. It is built on questions that more than one of the four interpretations
% answers: what the signs in the heavens are; why the disciples lift their
% heads; how Christ's coming is already under way and why the hour stays
% hidden; how a sleeper obeys the command to rise; what is near and why the
% Church still asks for Christ; what the soul of Psalm 24 is lifted from and
% toward, and what its waiting rests on; what the Communion's earth yields;
% and where the cleansing the orations ask for comes from. Under
% each, the authorities of the different interpretations are set side by side
% at the text they share. A short comparison of the four senses closes it.
% Every claim, attribution and qualification is the expansive study's; the
% connecting prose is written for this companion.

\runninghead{Detailed Commentary}
\section{The Propers: Detailed Commentary}
\zlabel{triptych:concise:commentary:start}

\subsection{The signs in the sun, the moon and the stars}

Luke says only that there will be signs in the sun, the moon and the stars;
Matthew's account of the same discourse says that the sun will be darkened,
the moon give no light and the stars fall (Mt 24:29). The Fathers divide on
whether these are portents still to come in the sky or figures of something
nearer.

Gregory takes them as real. The Lord, he says, foretells the evils that will
follow upon an aging world \latin{ut nos ab ejus amore compescat}, ``to
restrain us from loving it,'' so that those who would not fear God in quiet
may fear his judgment when worn down by blows (\work{Hom.\ in Ev.}~I.1.1). He
counts what his hearers had seen, nation rising against nation, earthquakes,
pestilence, and ``fiery battle lines in the sky'' before Italy fell to the
sword of the heathen, and argues from it: ``the showing of things past is the
certainty of things to follow.'' Later he likens the world to a body ``weighed
down by its very old age'' (I.1.5). Cyril of Alexandria reads the withering
as literal. Christ ``transfers His words'' from Jerusalem ``to the time of the
consummation''; creation ``begins, so to speak, to be changed,'' and ``the
unendurable fear of those things that are coming will suffice for the
destruction of many'' (\work{Comm.\ in Luc.}, sermon 139).

Augustine raises the strongest objection to this way of reasoning, in a
letter to a bishop, Hesychius, who had written that men were plainly
withering for fear in their own day (Letter 199.36). He replies first from
history. The histories of the nations record wonders before Christ came, two
suns among them; no one has seen the sun darkened as it was ``when the Light
of the world hung on the tree,'' at the full moon of the Passover, when it
cannot be eclipsed; and under Gallienus, when barbarians overran the
provinces, Christians could have thought the end was near (199.34--35). He
replies next from Scripture. Destruction comes ``when they shall say, Peace
and security'' (1~Thess 5:3), so the tribulation must fall on the Church and
not on those who trouble her; but the evils of his day fall alike on Christ's
people and the devil's, and the banquets and the dice go on: ``Is this
withering for fear, or rather soaking in lust?'' (199.36--37). The signs are
better understood of the Church, he concludes, lest the Lord seem to have
foretold as a great thing what befell the world before he came (199.39).

Augustine and Ambrose, who expounded the same verses of Luke, give the lights
different figures. For Augustine the three are one: \latin{Ecclesia est enim
sol et luna et stellae}, ``for the Church is sun and moon and stars,'' since
the Song of Songs calls her ``fair as the moon, bright as the sun'' (Cant
6:9). When persecutors rage and the wicked say ``Peace,'' the lights fail
together and the Church will not appear; the falling stars are ``many who
seemed to shine with grace'' and give way (199.39). For Ambrose the sun is
Christ: ``that heavenly sun is lessened or increased for me by my faith,'' as
the visible sun looks paler or brighter to each according to his sight. The
moon is the Church, which cannot borrow light from Christ's rays when the
vices of the flesh come between, and in persecution ``it is often love of this
life alone that shuts out the light of God'' (\work{Exp.\ in Luc.}~X.37). His
falling stars are the patriarchs and prophets, ``already shining with the
glory of the resurrection,'' lost to men's sight when persecution grows
bitter, ``for so the good are proved, and so the weak are exposed'' (X.38).
What the two share is the Church under persecution, when many fall from the
faith and her light is lost to sight, dimmed by the flesh for Ambrose and
hidden for Augustine. They differ on what the sun and the stars are.

The disagreement with Gregory is real, and both sides narrow it. Neither
Ambrose nor Augustine removes the last coming: Ambrose offers his figure of
the coming with \latin{Fortasse}, ``Perhaps'' (X.39), and Augustine excepts
what plainly belongs to the end (199.45). Gregory calls the calamities of his
time \latin{sequentis irae praecones}, ``heralds of the wrath to follow,'' as
unlike the last tribulation ``as far as the person of a herald is from the
power of a judge'' (I.1.6). Augustine's warning that the man who says the
Lord will come sooner ``speaks more desirably, but is dangerously deceived''
(199.54) still stands against Gregory's sense that the end was close. What
Gregory draws from the signs does not depend on it: they wean the heart from
a passing world, and they teach fear of the judge while there is time to
amend.

\subsection{Why the disciples lift their heads}

The Gospel's own difficulty is how the same events can wither some men and
lift the heads of others. For Gregory the difference lies in what each one loves. Those who
would not hear Christ in his humility, he says, will see him in power and
majesty (I.1.2); then the words are turned to the comfort of the elect:
\latin{levate capita, id est exhilarate corda}, ``lift up your heads, that
is, gladden your hearts, because while the world of which you are not friends
is ending, the redemption you have sought draws near'' (I.1.3). The head
stands for the mind, which orders the thoughts as the head rules the members.
Whoever does not rejoice as the end approaches shows, by St James's rule, that
he is the world's friend (Jas 4:4). Then Gregory asks, \latin{Quid est vita
mortalis, nisi via?} ``What is mortal life but a road? And consider, my
brothers, what it is to grow weary in the labour of the road, and yet not to
want the road to end.'' The fig tree, he says, shows that the world is to be
trodden underfoot: \latin{fructus mundi ruina est. Ad hoc enim crescit, ut
cadat}, ``the fruit of the world is its ruin. For it grows in order to fall.''
The summer is the kingdom, ``because then the clouds of our mourning pass, and
the days of life shine with the brightness of the eternal sun'' (I.1.3).

Gregory explains why the disciple may rejoice; Cyril of Alexandria explains
what he looks up for. Christ comes ``not secretly nor obscurely, but as God
and Lord,'' and ``will renew creation, and refashion the nature of man to that
which it was at the beginning''; the redemption is the resurrection, when
``this earthly and infirm body shall put off corruption'' (sermon 139). St
Cyril of Jerusalem says the like of the sky, teaching from Matthew's parallel:
``the Lord rolleth up the heavens, not that He may destroy them, but that He
may raise them up again more beautiful'' (\work{Cat.}~15.3). St Bede takes the
heaven that passes in verse~33 for the air, so that heaven and earth ``will
both pass away and be,'' cleansed of their present form and kept in their
nature (\work{In Lucam} VI, on 21:33). Of that verse Gregory observes that
nothing in nature lasts longer than heaven and earth and nothing passes faster
than a word; yet Christ's ``speech, which passes, expresses judgments that
remain without change'' (I.1.4).

Three days before he preached, a whirlwind had uprooted trees, destroyed
houses and overturned churches, and many who went to bed well died in the
night (I.1.5). If the judge did that with ``the breath of the slightest
wind,'' Gregory asks, what will he do ``when he comes in person''? His last
counsel is to amend and weep for sin, because ``you will see the coming of the
eternal judge the more securely, the more you now forestall his severity by
fearing it'' (I.1.6).

\subsection{The coming already under way, and the hidden hour}

The second interpretation finds Christ's coming already in progress.
Augustine sorts the discourse by what it foretells: the fall of the city,
Christ's coming in his body the Church, and his coming at the end in the
Church's head (199.25--26). In the Church he does not cease to come, and he is
recognised coming in his own ``while his members are born daily'' (199.25).
The cloud of his coming may be the Church or his own body, and Augustine says
the choice is hard (199.41--42). He reads on to the fig tree: among the things
after which the kingdom is near is the Son of man seen in the cloud, so when
he is seen the kingdom is still only near. Perhaps, then, the whole passage
may be taken of the coming by which ``he comes daily in his body, which is the
Church,'' ``as in a great cloud, bearing fruit and growing in the whole
world,'' except what plainly belongs to the last coming (199.42, 45). Ambrose
places the coming in each person. With his \latin{Fortasse} he likens it to a
presence ``that surely comes about in each one, when they have received Christ
with all their affections,'' so that the powers of heaven are moved by an
increase of grace (X.39). His clouds are the prophets and saints, and Christ
came ``in a light cloud, made flesh of the Virgin'' (X.42).

Aquinas gives the Epistle a sense of the same shape. The day at hand that
Paul meant is the day of glory; it can also be the day of grace, ``which,
though it has already come in the course of time, is said to draw near to us
by faith and devotion'' (\work{Super Rom.}~13, lect.~3). Christ has come, and
he still comes to each one who believes.

On ``this generation shall not pass away till all things be fulfilled'' (Lk
21:32) the second interpretation turns to Chrysostom on Matthew. He reads the
discourse of real trials, and the generation as one that outlasts them: ``not
of the generation then living, but of that of the believers,'' as the psalm
speaks of ``the generation of them that seek the Lord'' (Ps 23:6). Over it
nothing shall prevail, ``not famine, not pestilence, not earthquake, nor the
tumults of wars'' (\work{Hom.\ in Mt.}~77). Chrysostom does not read the
lights as figures of the Church, but the people of believers he describes is
the Church that Ambrose and Augustine see tried. Bede gives two other senses,
the whole human race or the Jewish people (on 21:32).

The hour of the end stays hidden. ``In whatever state his own last day finds
each one,'' Augustine writes, ``in that state the last day of the world will
overtake him'' (199.2), so the watch Christ commands after the lesson is
daily. Cyril of Jerusalem tells his catechumens, ``as to the time let no one
be curious \dots\ nor on the other hand supinely slumber'' (15.4). The third
interpretation measures the nearness differently. For Chrysostom each day
brings resurrection and judgment closer; Augustine refuses to measure the
time, and holds the desire for the Lord's coming right whether he comes soon
or late (199.1).

\subsection{The hour to rise, and who sets the sleeper free}

Paul's \latin{hora est iam nos de somno súrgere} raises a practical question:
how does a sleeper obey an order to get up? Chrysostom hears in it the hour of
the end, ``near is the Resurrection, near the awful Judgment,'' and sees a
danger in converts, whose earnestness would ``wither down to nothing'' as time
went on. Paul reverses that decline: ``The nearer the prize is, the more wide
awake ought they to be for the contest, since even the racers do this''
(\work{Hom.\ in Rom.}~24, on v.~11). The night far spent becomes a household
before dawn: at the swallow's twittering ``we each of us awake our neighbor,''
and at daybreak ``we all set about the works of the day, dressing, and leaving
our dreams'' (on v.~12). The armour of light is ``not in fact war, but a
solemn dance and feast-day,'' and its wearer ``at once soldier and
bridegroom''; at the last Paul ``gives us the Lord Himself for a garment, the
King Himself'' (on v.~14).

Aquinas asks what the sleep is: not of nature, nor the rest of contemplation,
but the sleep of sin, as in ``Rise, thou that sleepest'' (Eph 5:14), or of
negligence, as in ``How long wilt thou sleep, O sluggard?'' (Prov 6:9). One
rises from the first by penance and from the second by care for good works
(lect.~3).

Augustine tells in the \work{Confessions} what it is to hear such a command
and not obey it. A man puts off shaking off sleep, he says, \latin{quamvis
surgendi tempus advenerit}, ``although the time for rising has come''; to the
call ``Awake, thou that sleepest'' he himself had only ``the drawling and
drowsy words: `Presently, lo, presently;' \dots\ But `presently, presently,'
had no present.'' His question was St Paul's (Rom 7:24) made his own,
\latin{miserum ergo me quis liberaret}, ``who then would set me free,
wretched as I am?'', and the answer was grace (\work{Confessions} VIII.5.12).
The \work{Confessions} tell a life and do not expound Romans; the third
interpretation takes from them why a command to wake is not enough by itself.
The Collect's verb of rousing, \latin{Excita}, has Christ's power for its
object, not the sleeper. Its last words,
\latin{te liberánte salvári}, ``that we may be saved, with thee setting us
free,'' ask Christ, with the same verb, for the deliverance Augustine asked
for.

William Durandus heard the waking in the chant: the cantor begins
\latin{Ad te levávi} low and lifts his voice, \latin{quod proprium est
excitantis}, ``which is the mark of one rousing others,'' as in the Epistle's
\latin{Hora est} (\work{Rationale} VI). Sicard of Cremona sets out the way and
the wandering: ``The way is to rise from sleep, to put on the arms of light,
to walk becomingly. The wandering is to idle in the night and do the works of
darkness'' (\work{Mitrale} V). In the first interpretation, Paul's night and
Gregory's road ask the same of the hearer, to hurry as the end draws near, for
``we ought to hasten to them as fast as we can'' (I.1.3); Gregory adds that
only the friend of the world grieves to see it come.

\subsection{What is near, and why the Church still asks for Christ}

Aquinas explains Paul's ``our salvation is nearer'' in two ways that are the
Apostle's own: eternal life is near in time, as the just come closer to the
end of their lives, and near by disposition, as charity grows and makes a
person readier for the wedding feast (Mt 25:10). Then he adds a third and says
whose it is: \latin{secundum quod Ecclesia haec verba adventus tempore legit},
``as the Church reads these words in the time of Advent,'' they speak of the
salvation of Christ's first coming, and Paul speaks ``in the person of all the
faithful who have been since the beginning of the world'' (lect.~3; Isa
56:1). In Advent one may say, as in the time before the Incarnation, ``now our
salvation, that is, Christ, is nearer.'' He allows one more sense,
\latin{tempus misericordiae}, ``the time of mercy,'' when a man begins to want
to leave his past sins. The Epistle's \latin{dies \dots\ appropinquávit} and
the Gospel's \latin{appropínquat redémptio} use one verb, and Sicard says the
Epistle is fitted to the Gospel's ``the kingdom of God is at hand'' by its
``our salvation is nearer''; Durandus repeats him (\work{Mitrale} V;
\work{Rationale} VI).

Between the lessons the Alleluia asks for that salvation as a gift. The
first interpretation hears in it the people Gregory describes, who have
sought redemption and are then told that it is near. The third, if Paul's
words bear Aquinas's last sense, hears in ``Shew us, O Lord, thy mercy'' a
request for the time of mercy he describes. The fourth begins from this
verse.

Augustine preached on Psalm 84 to a congregation that had sung it: \latin{Hoc
est quod cantavimus}, ``This is what we sang.'' He reads the salvation as
Christ and meets the objection that Christ has been given. ``He has already
given us his Christ; yet let us still say to him, Give us thy Christ,'' as we
ask for daily bread, and the bread is the one who came down from heaven (Jn
6:51). The gift is not complete: \latin{Dedit enim nobis Christum, sed
hominem: quem nobis dedit hominem, eum ipsum nobis daturus est Deum}, ``For he
gave us Christ, but as man: the same whom he gave us as man he will give us as
God.'' He gave him as man because men could receive him only so. The second
gift is the sight of God, which Augustine explains by adoption: ``We are so
already, but by faith; we are so in hope, not yet in reality,'' for ``it hath
not yet appeared what we shall be \dots\ we shall see him as he is'' (1~Jn
3:2). A poor wanderer told that he is a senator's son would rejoice, if the
messenger were no deceiver; St John, no deceiver, tells those who despair,
``You have a father, you have a homeland, you have an inheritance''
(\work{En.\ in Ps.}~84.9). The mercy the verse asks to be shown is, for
Augustine, God's persuading a man that whatever good he has is from him, so
that he is not proud.

Bellarmine reads the psalm as a prophecy of Christ. On verse~8 the prophet
``now asks for the coming of the Savior,'' and ``grant us thy salvation''
means ``Send us your Son for a Savior'' (\work{Explanatio}, Ps 84, on v.~8).
He and Augustine agree that the salvation is Christ given in the flesh;
Augustine looks further, to the same Christ given as God. If in Advent the
Church takes the place, as Aquinas says, of the time when the Incarnation was
drawing near, she stands where Bellarmine puts the psalmist, among those who
ask for the Saviour before he comes. The second interpretation takes one
clause of Augustine's sermon, that those who sing are sons of God by faith and
not yet in reality: in it the Alleluia is sung by a Church to whom Christ has
come and in whom he still comes, and who has not yet appeared as what she will
be.

The Collect's \latin{et veni}, ``and come,'' makes to Christ himself the
request the Alleluia makes of the Father, and Augustine's reason for still
singing the one serves for the other. The fourth interpretation also relates
Augustine's two gifts to two things the Mass looks toward: the
Postcommunion's coming feast of the birth of Christ given as man, and the
Gospel's Son of man in majesty, when ``we shall see him as he is.''
Augustine does not call his two gifts comings; the two comings he names, on
Psalm 24, are ``the one in mercy, the other in judgment'' (\work{En.\ in
Ps.}~24.10). Cyril of Jerusalem preaches both
comings plainly: the first ``unobserved, like rain on a fleece,'' in the
manger, the second when ``He covereth Himself with light as with a garment''
(15.1).

\subsection{The lifted soul and those who wait}

The Introit's psalmist lifts his soul to God, and three of the authors take
up from what and toward what it is lifted. Gregory, explaining the Gospel's
\latin{leváte}, writes that \latin{Levare itaque capita est mentes nostras ad
gaudia patriae coelestis erigere}, ``to lift up the head is to raise our
minds to the joys of the heavenly homeland'' (I.1.3). Aquinas, on the psalm
itself, reads the lifting as the preparation for prayer, which is, in St John
Damascene's words, ``the ascent of the mind to God''; the psalmist lifts his
soul ``by contemplation, looking on thy goodness, and by love,'' and prays
with a trust placed, as Paul says, ``through Christ, towards God'' (2~Cor
3:4--5; \work{Super Ps.}~24, n.~1). Both make the lifting a raising of the
mind. For Gregory it rises toward the homeland while the world it does not
love is ending, and the first interpretation takes the psalmist, lifting his
soul while enemies threaten to laugh, to be doing already what the Gospel
commands of disciples who see the world shaken. For Aquinas it rises to God
in prayer.

Gregory and Aquinas speak of where the soul goes; Augustine, of what it
leaves. He hears in the psalm ``Christ \dots\ in the person of the Church,''
lifting the soul ``that was trodden down
on the earth with carnal longings,'' and he takes the enemies who would laugh
for ``secret suggestions'' that flatter a man to his fall with ``Well done,
well done'' (\work{En.\ in Ps.}~24.1--3). The second interpretation reads the
Gospel's signs of the Church's age of trial, and so hears the psalm as her
voice in that age. The third finds a soul pressed down in the same way in
Augustine's account of his delay: \latin{sarcina saeculi, velut somno
adsolet, dulciter premebar}, ``I was sweetly weighed down by the world's baggage, as one is by sleep''
(VIII.5.12). When the antiphon returns at the offering, after Paul's command
to rise, its lifted soul is the opposite of that weight.

The Introit's verse and the Gradual ask God to show his ways, and Aquinas and
Bellarmine both take the ways for God's law. Aquinas calls them the
commandments and says of them, \latin{Per has enim vias venit dominus ad nos,
maxime per praeceptum caritatis}, ``for by these ways the Lord comes to us,
above all by the commandment of charity'' (\work{Super Ps.}~24, n.~4). The
second interpretation finds there a coming in the present, of the kind
Ambrose and Augustine find in the Gospel. Bellarmine reads the request as a
prayer to be taught ``not speculatively, but practically,'' by grace that
moves the will (\work{Explanatio}, Ps 24, on v.~4). Aquinas's ways bring
the Lord to the one who keeps them; in Bellarmine's account the one who prays
asks for the will to keep them. The third interpretation, for which a command to wake is not enough by itself, follows
him.

The chants sing that those \latin{qui te exspéctant}, ``who wait for thee,''
are not confounded, and the Gospel uses the noun of the same verb for men who
wither \latin{præ timóre et exspectatióne}, ``for fear and expectation.'' The
first interpretation distinguishes the two by what is awaited. Those men
expect what will come upon the world, whose ruin, Gregory has said, grieves
only those who love it; those who wait for God wait for him. In the fourth
interpretation Bellarmine explains the psalm's waiting: to wait on the Lord
``means to expect him in the certain hope of assistance'' (on vv.~2--3). That
interpretation also names the one awaited in this Mass, the Christ whom the
Alleluia and the Collect ask for.

\subsection{What the earth yields}

Rupert of Deutz, expounding the Communion in his commentary on the Mass,
reads it of the Incarnation: \latin{terra nostra, id est, uterus Virginis
concipiens, dedit fructum suum}, ``our earth, that is, the womb of the Virgin
conceiving, gave its fruit'' (\work{De div.\ off.}~III.2). Sicard repeats it
(\work{Mitrale} V). Bellarmine gives it its fullest form: the truth springs
out of the earth ``not in the manner of the seed that we sow and cultivate,
but in the manner of the natural flowers that grow spontaneously''; the
Lord's goodness is the Holy Spirit, ``who will overshadow a virgin, and thus
our land, which was never ploughed nor sown \dots\ will yield her fruit'' (on
v.~13).

Augustine finds in the verse the grace given to the believer. He takes
``Truth is sprung out of the earth'' of the confession of sins, and ``justice
hath looked down from heaven'' of the justification God gives to the one who
confesses. The sweetness is ``the sweetness of doing justice'': the drunkard
delights in sobriety, the thief looks for someone to give to, and ``you run
to the church, who used to run to the theatre.'' Then he speaks to his
hearers as a field furrowed by ``the plough of confession,'' where he has sown
the word: ``And unless God sends rain, what use is what is sown? \dots\ Let
God's rain come'' (84.15). Sicard gives a second explanation close to
Augustine's, the goodness a good affection, the earth our flesh, its fruit
good works, and Durandus copies him.

The fourth interpretation keeps both explanations, and they agree that the
earth bears nothing of itself: Bellarmine's land was ``never ploughed nor
sown,'' and Augustine's field needs God's rain. They differ on the fruit. For
Rupert and Bellarmine it is Christ, born once of the Virgin, a sense that
suits a Mass preparing for Christmas; for Augustine it is the good works that
God's sweetness makes grow wherever his rain falls, a sense that suits the
moment of Communion, when the faithful receive what they asked for.

The second and third interpretations take the fruit for good works, with
Augustine and Sicard's second explanation, and leave the Virgin aside. The second looks to the rain. Augustine's cloud of the
daily coming bears fruit in the whole world, and in the same letter he takes the clouds of
the coming to be those God can command not to rain (Isa 5:6; 199.25); if the
rain he prays for in the sermon falls from those clouds, the Communion's
harvest is the fruit of Christ's coming in his Church. In the third the
earth is our flesh, as Sicard has it, the flesh Paul would keep from
revelling and wantonness, and its fruit is the good works by which, Aquinas
said, a man rises from the sleep of negligence. The first interpretation
sets the harvest against a different fruit, the world's, which Gregory called
its ruin because the world grows in order to fall. The Communion, it notes,
promises another harvest and begins with the Lord's gift, ``The Lord will
give goodness''; the third interpretation notes the same order, the Lord's
gift before the earth's yield.

\subsection{The Collect, the Secret and the Postcommunion}

Schuster paraphrases the Collect's petition as a request that our Lord ``save
us by his power from the dangers to which we are exposed through our sins''
(p.~321), and the Secret asks the holy gifts to cleanse us \latin{poténti
virtúte}, by their mighty power. The interpretations differ on where the
cleansing from sin comes from. The first sets both prayers beside Gregory.
The danger the Collect names lies in the one who prays, as does the danger
Gregory has in view when he tells his hearers to amend; but Gregory puts the
remedy in tears and changed ways, and the Secret asks the cleansing of the
gifts themselves. The third sets the Secret beside Paul, who makes the
casting off of the works of darkness the hearers' own task, and the Secret
asks God for a cleansing like it. The first and third interpretations thus
make the same contrast, between what Gregory and Paul tell the sinner to do
and what the Secret asks of God. In the second interpretation the cleansing
removes an obstacle: Ambrose's moon loses the sun's light when the vices of
the flesh stand in the way, and the Secret asks for the cleansing that would
take them away. The fourth is concerned with what the cleansing prepares
for: the
Author whom Schuster names in the Secret would be the Christ whom the
Alleluia asked God to give.

With Schuster, the third interpretation takes the Postcommunion's coming
solemnities of our restoration for Christmas, the feast of the first coming,
as Aquinas says the Church reads Paul in Advent. The first interpretation
keeps Christmas and adds the restoration Cyril of Alexandria describes, which
he places at Christ's return: creation renewed and man's nature made as it
was at the beginning. The prayer is said after a Gospel of that return, and
the birth it prepares to keep would then begin the restoration that the
second coming completes. The prayer asks for mercy in order to go before
that feast. If that mercy is the one Augustine described on the Alleluia's
verse, the persuasion that whatever good a man has is from God, then in the
fourth interpretation the communicant asks to receive that. The second
attends instead to the place, \latin{in médio templi tui}, in the midst of
the temple, which would be one place of Christ's daily coming in his body. In that interpretation the Collect's \latin{veni},
prayed after the Church's psalm, may likewise ask for the daily coming as
well as for the last.

\subsection{The four senses compared}

The interpretations differ most in their allegorical senses. In the first the
allegory is the coming itself: Christ, hidden and unheard at his first coming,
comes openly in power, and those who love him find him as the world they did
not love passes; its anagogy is the body's resurrection and creation renewed.
In the second it is the Church, whose failing lights Ambrose and Augustine
assign each in his own way, and whose cloud is the Church or Christ's body
(Augustine), or the prophets and saints, and the flesh taken of the Virgin
(Ambrose); its anagogy is the coming in the head that completes the coming in
the body. The third has the least allegory, Aquinas's Advent sense of the
Epistle and Chrysostom's Christ put on as a garment, and the most moral
sense; its anagogy is the day of glory, nearer every day, which the one who
has put on Christ goes out to meet as soldier and bridegroom (Chrysostom;
Aquinas on the Apostle's meaning). The fourth's allegory is Christ himself, the salvation asked for and, for Rupert
and Bellarmine, the fruit of the untouched earth; its anagogy is the sight of
him as he is.

Their moral senses set different things before the hearer: the world's ruin,
which he must not love, and the judge he must fear now (Gregory); the state of
faith in himself and in the Church, with a daily watch because the hour is
hidden (Ambrose, Augustine); his own sleep, from which he rises by penance and
good works and from which Christ must set him free (Aquinas, Chrysostom, the
Collect); and his pride, which God's mercy cures by teaching that all his good
is from God (Augustine). On present events the first two differ: Gregory
reads the disasters of his time as heralds of the judge, though as unlike him
as a herald is unlike a judge, and Augustine warns that dating the end from
such things is a dangerous error. The third
and fourth both look toward Christmas, the third through the hearer, who must
wake, the fourth through the gift, Christ asked for though already given. The
fourth alone sets the two explanations of the Communion's earth against each
other; the second and third read its fruit as good works without the dispute.
